% Reliability diagram on test held-out responses; artifacts/evaluation/calibration_test.csv
\begin{tikzpicture}
\begin{axis}[
  width=\linewidth, height=\figH, xmin=0, xmax=1, ymin=0, ymax=1,
  xlabel={rata-rata prediksi P(benar)}, ylabel={proporsi benar teramati},
  grid=major, legend pos=north west, legend cell align=left,
  xtick={0,0.2,...,1}, ytick={0,0.2,...,1},
  ticklabel style={/pgf/number format/fixed, /pgf/number format/precision=1},
]
  \addplot[dashed, midgray, domain=0:1] {x};
  \addlegendentry{kalibrasi sempurna}
  \addplot[navy, thick, mark=*, mark size=1.4pt] coordinates {
    (0.0754,0.0000) (0.1721,0.1808) (0.2588,0.2608) (0.3445,0.3536) (0.4470,0.4575)
    (0.5467,0.5436) (0.6497,0.6354) (0.7476,0.7628) (0.8541,0.8502) (0.9414,0.9465)};
  \addlegendentry{P-DINA (ECE 0,008)}
  \addplot[tealc, thick, mark=square*, mark size=1.3pt] coordinates {
    (0.0780,0.1940) (0.1572,0.2362) (0.2525,0.2615) (0.3517,0.3330) (0.4499,0.3950)
    (0.5500,0.5172) (0.6492,0.6279) (0.7500,0.7755) (0.8476,0.9073) (0.9346,0.9779)};
  \addlegendentry{Q-logistic (ECE 0,036)}
\end{axis}
\end{tikzpicture}
